\chapter{Konvergenz reeller Folgen}
\label{b21guide:ch:limits}

Von hier an sind $a,b,u,v$ reelle Folgen und $L,K,c$ reelle
Zahlen, sofern nichts anderes angegeben wird.
Es gelten die übliche reelle Arithmetik und der reelle Betrag.
Ein Grenzwert ohne weiteren Zusatz ist stets endlich und reell.

\section{Definition und Grundkriterien}

\begin{definition}[Grenzwert, Konvergenz und Divergenz]
\label{b21guide:def:limit}
Die Schreibweise $a_n\to L$ oder
$\lim_{n\to\infty}a_n=L$ bedeutet
\[
 \forall\varepsilon>0\;\exists N\in\mathbb N\;
 \forall n\geq N:\ |a_n-L|<\varepsilon.
\]
Eine Folge heißt \emph{konvergent}, wenn ein solches $L\in\mathbb R$
existiert, sonst \emph{divergent}.
Eine gegen $0$ konvergente Folge heißt \emph{Nullfolge}.
\end{definition}

\begin{theorem}[Eindeutigkeit und konstante Folgen]
Eine reelle Folge hat höchstens einen Grenzwert.
Eine konstante Folge $a_n=c$ konvergiert gegen $c$.
Dasselbe gilt, wenn $a_n=c$ nur schließlich erfüllt ist.
\end{theorem}

\begin{theorem}[Abzählbares Prüfkriterium und Negation]
\label{b21guide:thm:epsilon-test}
Es gilt $a_n\to L$ genau dann, wenn für jedes ganze $k\geq1$
schließlich $|a_n-L|<1/k$ gilt.
Die Negation eines vorgegebenen Grenzwerts lautet
\[
 a_n\not\to L
 \quad\Longleftrightarrow\quad
 \exists\varepsilon_0>0\;\forall N\in\mathbb N\;
 \exists n\geq N:\ |a_n-L|\geq\varepsilon_0.
\]
\end{theorem}

\begin{theorem}[Rückführung auf Nullfolgen]
\[
 a_n\to L
 \quad\Longleftrightarrow\quad a_n-L\to0
 \quad\Longleftrightarrow\quad |a_n-L|\to0.
\]
\end{theorem}

\begin{theorem}[Endliche Änderungen und Restfolgen]
\label{b21guide:thm:tail-limit}
Gilt $a_n=b_n$ schließlich, so ist
$a_n\to L$ äquivalent zu $b_n\to L$.
Für jedes feste $r\in\mathbb N$ gilt
\[
 a_n\to L\quad\Longleftrightarrow\quad a_{n+r}\to L.
\]
Eine Änderung, Hinzufügung oder Entfernung endlich vieler
Anfangsglieder ändert somit den Grenzwert nicht.
\end{theorem}

\section{Beschränktheit und Vergleiche}

\begin{definition}[Beschränkte Folge]
\label{b21guide:def:bounded}
Eine reelle Folge heißt \emph{beschränkt}, wenn
\[
 \exists C\geq0\;\forall n\in\mathbb N:\ |a_n|\leq C.
\]
Sie heißt \emph{nach oben beschränkt}, wenn $a_n\leq C$ für alle
$n$ und eine reelle Zahl $C$ gilt; \emph{nach unten beschränkt}
bedeutet entsprechend $c\leq a_n$.
\end{definition}

\begin{theorem}[Konvergenz erzwingt Beschränktheit]
\label{b21guide:thm:bounded}
Jede konvergente reelle Folge ist beschränkt.
Eine reelle Folge ist bereits dann beschränkt, wenn irgendeine
ihrer Restfolgen beschränkt ist.
Beschränktheit allein ist kein Konvergenzkriterium.
\end{theorem}

\begin{theorem}[Schließliche Trennung vom Grenzwert]
\label{b21guide:thm:separation}
Aus $a_n\to L$ folgen
\[
 \begin{aligned}
 c<L&\ \Longrightarrow\ a_n>c\text{ schließlich},\\
 L<c&\ \Longrightarrow\ a_n<c\text{ schließlich}.
 \end{aligned}
\]
Ist $L\ne0$, so hat $a_n$ schließlich dasselbe Vorzeichen wie $L$;
insbesondere gilt schließlich $|a_n|\geq |L|/2>0$.
\end{theorem}

\begin{theorem}[Vergleiche beim Grenzübergang]
\label{b21guide:thm:comparison}
Seien $a_n\to L$ und $b_n\to K$.
Gilt $a_n\leq b_n$ schließlich, so ist $L\leq K$.
Auch aus $a_n<b_n$ schließlich folgt im allgemeinen nur $L\leq K$.
Umgekehrt erzwingt $L<K$ die schließliche Ungleichung $a_n<b_n$.
\end{theorem}

\begin{theorem}[Majorantenkriterium für Nullfolgen]
\label{b21guide:thm:majorant}
Gelte $b_n\to0$ und schließlich $|a_n|\leq b_n$.
Dann gilt $a_n\to0$.
Allgemeiner folgt aus $b_n\to0$ und
$|a_n-L|\leq b_n$ schließlich die Konvergenz $a_n\to L$.
\end{theorem}

\begin{theorem}[Einschließungsprinzip]
\label{b21guide:thm:squeeze}
Gelte $u_n\leq a_n\leq v_n$ schließlich sowie
$u_n\to L$ und $v_n\to L$. Dann ist $a_n\to L$.
\end{theorem}

\begin{theorem}[Verschwindender Unterschied]
\label{b21guide:thm:perturbation}
Ist $a_n-b_n\to0$, so gilt für jedes $L\in\mathbb R$
\[
 a_n\to L\quad\Longleftrightarrow\quad b_n\to L.
\]
\end{theorem}

\section{Grenzwerterhaltende Indexfolgen}

\begin{definition}[Gegen unendlich laufende Indizes]
\label{b21guide:def:escaping-indices}
Für eine Folge $(r_n)$ natürlicher Zahlen bedeutet $r_n\to+\infty$:
\[
 \forall K\in\mathbb N\;\exists N\in\mathbb N\;
 \forall n\geq N:\ r_n\geq K.
\]
Die Indexfolge muß weder monoton sein noch paarweise verschiedene Werte besitzen.
\end{definition}

\begin{theorem}[Grenzwerterhaltung bei Umindizierung]
\label{b21guide:thm:reindex-limit}
Für eine Folge $(r_n)$ natürlicher Zahlen sind äquivalent:
\begin{enumerate}
 \item $r_n\to+\infty$.
 \item Für jede reelle Folge $(a_n)$ und jedes $L\in\mathbb R$
 gilt: Aus $a_n\to L$ folgt $a_{r_n}\to L$.
\end{enumerate}
Damit genügt zum Erhalten eines Grenzwerts das schließliche
Verlassen jedes endlichen Indexabschnitts. Bloße Unbeschränktheit
der Indexfolge genügt nicht.
\end{theorem}
